% !TEX root = ../main.tex
\chapter{Analyse Exploratoire des Donn\'ees (EDA)}
\label{ch:eda}

\begin{intuition}
L'analyse exploratoire des donn\'ees (EDA) est une philosophie d'analyse initi\'ee par John Tukey dans les ann\'ees 1970. Plut\^ot que de tester des hypoth\`eses pr\'ed\'efinies, l'EDA invite \`a <<~laisser parler les donn\'ees ~>> en les explorant syst\'ematiquement par des r\'esum\'es num\'eriques et des visualisations.
\end{intuition}

% ============================================================
\section{Qu'est-ce que l'EDA ?}
\index{EDA}\index{analyse exploratoire}

\begin{definition}[Analyse exploratoire des donn\'ees]
L'EDA\index{EDA} est une approche d'analyse qui utilise des techniques visuelles et quantitatives pour comprendre la structure, les anomalies et les relations dans un jeu de donn\'ees \emph{avant} toute mod\'elisation formelle.
\end{definition}

Les objectifs principaux de l'EDA sont :
\begin{itemize}
    \item D\'ecouvrir la \textbf{structure} des donn\'ees (dimensions, types, valeurs manquantes).
    \item Identifier les \textbf{distributions} et les \textbf{tendances centrales}.
    \item Rep\'erer les \textbf{anomalies} et les \textbf{valeurs aberrantes}.
    \item Explorer les \textbf{relations} entre variables.
    \item Formuler des \textbf{hypoth\`eses} \`a tester ult\'erieurement.
\end{itemize}

% ============================================================
\section{Processus syst\'ematique d'EDA}

\begin{center}
\begin{tikzpicture}[
    node distance=1.2cm,
    block/.style={rectangle, draw=steelblue!80, fill=steelblue!10, rounded corners, text width=3.8cm, align=center, minimum height=1cm, font=\small},
    arrow/.style={->, thick, >=stealth, color=gray!70}
]
\node[align=center,block] (load) {1. Charger les donn\'ees\\{\scriptsize \py{pd.read\_csv}, \py{sns.load\_dataset}}};
\node[align=center,block, below=of load] (shape) {2. Dimensions et types\\{\scriptsize \py{shape}, \py{dtypes}, \py{info()}}};
\node[align=center,block, below=of shape] (missing) {3. Valeurs manquantes\\{\scriptsize \py{isnull().sum()}}};
\node[align=center,block, below=of missing] (desc) {4. Statistiques descriptives\\{\scriptsize \py{describe()}, \py{value\_counts()}}};
\node[align=center,block, right=2.5cm of load] (uni) {5. Analyse univari\'ee\\{\scriptsize histogrammes, boxplots}};
\node[align=center,block, below=of uni] (bi) {6. Analyse bivari\'ee\\{\scriptsize scatter, corr\'elation}};
\node[align=center,block, below=of bi] (multi) {7. Analyse multivari\'ee\\{\scriptsize pairplot, heatmap}};
\node[block, below=of multi] (concl) {8. Conclusions et hypoth\`eses};

\draw[arrow] (load) -- (shape);
\draw[arrow] (shape) -- (missing);
\draw[arrow] (missing) -- (desc);
\draw[arrow] (desc.east) -- ++(1.25,0) |- (uni.west);
\draw[arrow] (uni) -- (bi);
\draw[arrow] (bi) -- (multi);
\draw[arrow] (multi) -- (concl);
\end{tikzpicture}
\end{center}

% ============================================================
\section{\'Etude de cas : le jeu Titanic}
\index{Titanic}

\subsection{Chargement et premi\`ere inspection}

\begin{codeblock}[title=\textbf{Python}]
\begin{minted}{python}
import pandas as pd
import numpy as np
import seaborn as sns
import matplotlib.pyplot as plt

titanic = sns.load_dataset('titanic')
print(f"Dimensions : {titanic.shape}")
print(titanic.dtypes)
\end{minted}
\end{codeblock}

\begin{codeoutput}
\begin{minted}{text}
Dimensions : (891, 15)
survived         int64
pclass           int64
sex             object
age            float64
sibsp            int64
parch            int64
fare           float64
embarked        object
class         category
who             object
adult_male        bool
deck          category
embark_town     object
alive           object
alone             bool
\end{minted}
\end{codeoutput}

\subsection{Valeurs manquantes}

\begin{codeblock}[title=\textbf{Python}]
\begin{minted}{python}
missing = titanic.isnull().sum()
missing_pct = (missing / len(titanic) * 100).round(1)
missing_df = pd.DataFrame({'Manquantes': missing,
                           'Pourcentage': missing_pct})
print(missing_df[missing_df['Manquantes'] > 0])
\end{minted}
\end{codeblock}

\begin{codeoutput}
\begin{minted}{text}
             Manquantes  Pourcentage
age                 177         19.9
embarked              2          0.2
deck                688         77.2
embark_town           2          0.2
\end{minted}
\end{codeoutput}

\begin{remarque}
La variable \py{deck} a 77\,\% de valeurs manquantes : elle sera probablement inutilisable sans imputation lourde. La variable \py{age} a environ 20\,\% de manquants, ce qui est g\'erable.
\end{remarque}

\subsection{Statistiques descriptives}

\begin{codeblock}[title=\textbf{Python}]
\begin{minted}{python}
print(titanic.describe())
print("\n--- Variables catégorielles ---")
print(titanic[['sex', 'embarked', 'class']].describe())
\end{minted}
\end{codeblock}

\begin{codeoutput}
\begin{minted}{text}
         survived     pclass        age      sibsp      parch       fare
count  891.000000  891.000000  714.00000  891.00000  891.00000  891.00000
mean     0.383838    2.308642   29.69912    0.52301    0.38159   32.20421
std      0.486592    0.836071   14.52650    1.10274    0.80606   49.69343
min      0.000000    1.000000    0.42000    0.00000    0.00000    0.00000
25%      0.000000    2.000000   20.12500    0.00000    0.00000    7.91040
50%      0.000000    3.000000   28.00000    0.00000    0.00000   14.45420
75%      1.000000    3.000000   38.00000    1.00000    0.00000   31.00000
max      1.000000    3.000000   80.00000    8.00000    6.00000  512.32920

--- Variables catégorielles ---
          sex embarked class
count     891      889   891
unique      2        3     3
top      male        S Third
freq      577      644   491
\end{minted}
\end{codeoutput}

% ============================================================
\section{Analyse univari\'ee}
\index{analyse univari\'ee}

\begin{codeblock}[title=\textbf{Python}]
\begin{minted}{python}
fig, axes = plt.subplots(2, 2, figsize=(12, 9))

sns.histplot(titanic['age'].dropna(), bins=30, kde=True,
             ax=axes[0,0], color='steelblue')
axes[0,0].set_title('Distribution de l\'\\^age')

sns.countplot(data=titanic, x='pclass', ax=axes[0,1],
              palette='viridis')
axes[0,1].set_title('R\'epartition par classe')

sns.countplot(data=titanic, x='survived', ax=axes[1,0],
              palette='Set2')
axes[1,0].set_xticklabels(['Décédé', 'Survivant'])
axes[1,0].set_title('Survie')

sns.histplot(titanic['fare'], bins=40, ax=axes[1,1],
             color='coral')
axes[1,1].set_title('Distribution du tarif')

plt.tight_layout()
plt.show()
\end{minted}
\end{codeblock}

\begin{codeoutput}
Quatre graphiques : l'\^age suit approximativement une distribution normale centr\'ee vers 29 ans ; la troisi\`eme classe est la plus repr\'esent\'ee (491) ; 62\,\% des passagers n'ont pas surv\'ecu ; le tarif est tr\`es asym\'etrique avec une longue queue droite.
\end{codeoutput}

% ============================================================
\section{Analyse bivari\'ee}
\index{analyse bivari\'ee}

\subsection{Survie par sexe et par classe}

\begin{codeblock}[title=\textbf{Python}]
\begin{minted}{python}
fig, axes = plt.subplots(1, 3, figsize=(15, 5))

# Survie par sexe
sns.barplot(data=titanic, x='sex', y='survived',
            ax=axes[0], palette='coolwarm', ci=95)
axes[0].set_title('Taux de survie par sexe')
axes[0].set_ylabel('Taux de survie')

# Survie par classe
sns.barplot(data=titanic, x='pclass', y='survived',
            ax=axes[1], palette='viridis', ci=95)
axes[1].set_title('Taux de survie par classe')

# Survie par sexe et classe
sns.barplot(data=titanic, x='pclass', y='survived',
            hue='sex', ax=axes[2], palette='Set1', ci=95)
axes[2].set_title('Survie par classe et sexe')

plt.tight_layout()
plt.show()
\end{minted}
\end{codeblock}

\begin{codeoutput}
Les femmes ont un taux de survie de 74\,\% contre 19\,\% pour les hommes. La 1\`ere classe a 63\,\% de survie, la 2\`eme 47\,\% et la 3\`eme 24\,\%. La combinaison montre que les femmes de 1\`ere classe ont surv\'ecu \`a pr\`es de 97\,\%.
\end{codeoutput}

\subsection{Corr\'elations et table de contingence}
\index{table de contingence}

\begin{codeblock}[title=\textbf{Python}]
\begin{minted}{python}
# Table de contingence
ct = pd.crosstab(titanic['sex'], titanic['survived'],
                 margins=True)
ct.columns = ['Décédé', 'Survivant', 'Total']
print(ct)

# Corrélation numérique
cols_num = ['survived', 'pclass', 'age', 'sibsp', 'parch', 'fare']
corr = titanic[cols_num].corr()
print("\nMatrice de corrélation :")
print(corr.round(2))
\end{minted}
\end{codeblock}

\begin{codeoutput}
\begin{minted}{text}
         Décédé  Survivant  Total
sex
female       81        233    314
male        468        109    577
Total       549        342    891

Matrice de corrélation :
          survived  pclass    age  sibsp  parch   fare
survived      1.00   -0.34  -0.08  -0.04   0.08   0.26
pclass       -0.34    1.00  -0.37   0.08   0.02  -0.55
age          -0.08   -0.37   1.00  -0.31  -0.19   0.10
fare          0.26   -0.55   0.10  -0.24   0.22   1.00
\end{minted}
\end{codeoutput}

% ============================================================
\section{Analyse multivari\'ee}
\index{analyse multivari\'ee}

\begin{codeblock}[title=\textbf{Python}]
\begin{minted}{python}
fig, ax = plt.subplots(figsize=(8, 6))
sns.heatmap(corr, annot=True, fmt='.2f', cmap='RdBu_r',
            center=0, ax=ax, square=True,
            linewidths=0.5)
ax.set_title('Heatmap des corr\'elations -- Titanic')
plt.tight_layout()
plt.show()
\end{minted}
\end{codeblock}

\begin{codeoutput}
Heatmap montrant les corr\'elations : la survie est n\'egativement corr\'el\'ee avec la classe ($-0.34$) et positivement avec le tarif ($0.26$). La classe et le tarif sont fortement anti-corr\'el\'es ($-0.55$).
\end{codeoutput}

% ============================================================
\section{D\'etection des valeurs aberrantes}
\index{valeurs aberrantes}\index{IQR}

\begin{definition}[M\'ethode IQR]
Une observation est consid\'er\'ee comme aberrante si elle se trouve en dehors de l'intervalle :
\[
\bigl[Q_1 - 1.5 \times \text{IQR},\; Q_3 + 1.5 \times \text{IQR}\bigr]
\]
o\`u $\text{IQR} = Q_3 - Q_1$ est l'\'ecart interquartile\index{\'ecart interquartile}.
\end{definition}

\begin{codeblock}[title=\textbf{Python}]
\begin{minted}{python}
def detecter_outliers_iqr(df, colonne):
    Q1 = df[colonne].quantile(0.25)
    Q3 = df[colonne].quantile(0.75)
    IQR = Q3 - Q1
    borne_inf = Q1 - 1.5 * IQR
    borne_sup = Q3 + 1.5 * IQR
    outliers = df[(df[colonne] < borne_inf) |
                  (df[colonne] > borne_sup)]
    return outliers, borne_inf, borne_sup

outliers_fare, b_inf, b_sup = detecter_outliers_iqr(titanic, 'fare')
print(f"Bornes : [{b_inf:.2f}, {b_sup:.2f}]")
print(f"Nombre d'outliers (fare) : {len(outliers_fare)}")
print(f"Tarif max : {titanic['fare'].max():.2f}")
\end{minted}
\end{codeblock}

\begin{codeoutput}
\begin{minted}{text}
Bornes : [-26.72, 65.63]
Nombre d'outliers (fare) : 116
Tarif max : 512.33
\end{minted}
\end{codeoutput}

\begin{remarque}
116 passagers (13\,\%) ont un tarif consid\'er\'e comme aberrant par la m\'ethode IQR. Ce sont principalement des passagers de premi\`ere classe. Il ne faut pas les supprimer automatiquement : ce sont des valeurs l\'egitimes qui refl\`etent la structure des donn\'ees.
\end{remarque}

% ============================================================
\section{Exercices}

\begin{exercice}[$\star$]
Chargez le jeu \py{titanic} et affichez les 5 premi\`eres lignes, les dimensions, les types de colonnes et le nombre de valeurs manquantes par colonne.
\end{exercice}

\begin{exercice}[$\star$]
Calculez le taux de survie global, puis le taux de survie par port d'embarquement (\py{embarked}). Quelle est la lettre du port avec le meilleur taux de survie ?
\end{exercice}

\begin{exercice}[$\star\star$]
Cr\'eez un graphique avec \py{sns.catplot} de type \py{'bar'} montrant le taux de survie par classe et par sexe. Ajoutez un titre informatif. Interpr\'etez les r\'esultats en termes de priorit\'e d'\'evacuation.
\end{exercice}

\begin{exercice}[$\star\star$]
Cr\'eez un histogramme de l'\^age des passagers, s\'epar\'e par survie (\py{hue='survived'}). Les survivants sont-ils en moyenne plus jeunes ? V\'erifiez avec \py{groupby('survived')['age'].mean()}.
\end{exercice}

\begin{exercice}[$\star\star\star$]
Appliquez la m\'ethode IQR pour d\'etecter les outliers dans la variable \py{age}. Combien sont-ils ? Cr\'eez un boxplot annot\'e montrant les bornes IQR et les outliers. Discutez de la pertinence de supprimer ces observations.
\end{exercice}

\begin{exercice}[$\star\star\star$]
R\'ealisez une EDA compl\`ete du jeu \py{penguins} (\py{sns.load\_dataset('penguins')}). Votre analyse doit inclure : (a) inspection g\'en\'erale, (b) traitement des valeurs manquantes, (c) distributions univari\'ees, (d) analyse bivari\'ee par esp\`ece, (e) heatmap de corr\'elation. Concluez en formulant 3 hypoth\`eses testables.
\end{exercice}

% ============================================================
\begin{formules}
\textbf{R\'esum\'e du chapitre -- Analyse exploratoire des donn\'ees}
\begin{itemize}
    \item \textbf{Inspection} : \py{df.shape}, \py{df.dtypes}, \py{df.info()}, \py{df.head()}.
    \item \textbf{Manquants} : \py{df.isnull().sum()}, \py{df.isnull().mean() * 100}.
    \item \textbf{Descriptif} : \py{df.describe()}, \py{df['col'].value\_counts()}.
    \item \textbf{Univari\'e} : \py{sns.histplot()}, \py{sns.countplot()}, \py{sns.boxplot()}.
    \item \textbf{Bivari\'e} : \py{sns.barplot(x, y)}, \py{sns.scatterplot()}, \py{pd.crosstab()}.
    \item \textbf{Multivari\'e} : \py{df.corr()}, \py{sns.heatmap()}, \py{sns.pairplot()}.
    \item \textbf{Outliers (IQR)} : aberrant si $x < Q_1 - 1.5 \cdot \text{IQR}$ ou $x > Q_3 + 1.5 \cdot \text{IQR}$.
    \item \textbf{Philosophie Tukey} : explorer d'abord, mod\'eliser ensuite.
\end{itemize}
\end{formules}
